\chapter*{Appendix A: The Transcription Standard}

Every reference transcript in this thesis was produced against the written standard
summarised here, version 1.0. The full document was given to the transcription team
before any file was written, and a validation script checks every file against it.

\section*{A.1 The segment rule}

The single most important rule is that no segment may be longer than 30 seconds,
because that is the length of audio the speech model sees at once. Anything longer is
cut by software that does not know where sentences end, which produces audio that does
not match its text. The target is 10 to 25 seconds per segment, and a break is always
taken at a natural pause, never inside a word or a phrase.

The nine transcripts produced before the standard existed do not follow this rule.
Their segments run to a median of 79 seconds and up to 225 seconds, and 98 per cent are
over 30 seconds. They are subdivided at sentence ends by the data preparation step, and
the report says so wherever they are used.

\section*{A.2 File format}

One UTF-8 text file per video, named \texttt{BanglaASR<N>\_ground\_truth.txt}, with a
header of comment lines and then one line per segment.

\begin{verbatim}
# BanglaASR12 Ground Truth Transcription
# Topic: Python Functions
# Speaker ID: B
# Video duration: 14:22
# Transcriber: RH
# Style guide version: 1.0
# ==========================================================================
[0:00-0:16] Assalamualaikum shobai. Aj amra dekhbo je kivabe python e
            function likhte hoy.
[0:16-0:33] Function mane hocche ekta code block jeta amra baar baar call
            korte pari.
\end{verbatim}

The timestamp is exactly \texttt{[M:SS-M:SS]} with no space anywhere inside the
brackets. The end timestamp of one segment is the start of the next, so there are no
gaps and no overlaps. Lines beginning with \texttt{\#} are ignored by the software.

\section*{A.3 Write it exactly as spoken}

Fillers are kept, every one of them, every time. This includes \textit{um},
\textit{uh}, \textit{aaa}, \textit{mane}, \textit{je}, \textit{tahole},
\textit{okay} and \textit{thik ache}. Three reasons were given to the team and all
three are real. The model has to learn them, and in the existing files fillers are
about 5 per cent of tokens, at 671 out of 12,890. Removing them breaks the alignment
between the text and its time window. And a reference cleaner than reality makes every
error rate reported from it wrong in our favour.

The one thing that is removed is an abandoned restart, where the lecturer begins a
sentence, stops and starts again differently. Only the finished version is kept. If it
is unclear whether something is a restart or a filler, it is kept.

The lecturer is never corrected. Slips, repeats and self-corrections are transcribed as
they happen.

\section*{A.4 Spelling}

Banglish has no spelling standard, so the guide fixes one. If one transcriber writes
\textit{kore} and another writes \textit{koree} for the same word, the model sees two
unrelated words and learns both badly. The canonical list was derived from what the
team already wrote most often, so it mostly matches their instincts. Table
\ref{tab:spellinglist} gives a sample of it.

\begin{table}[H]
\centering
\caption{A sample of the canonical spelling list. The validator enforces these.}
\label{tab:spellinglist}
\small
\begin{tabular}{p{4.0cm}p{3.0cm}p{5.0cm}}
\toprule
\textbf{Meaning} & \textbf{Use this} & \textbf{Do not write} \\
\midrule
I, my & \texttt{ami}, \texttt{amar} & amr, aami \\
we & \texttt{amra} & aamra, amraa \\
this & \texttt{eta}, \texttt{ei} & etaa, eee \\
one, a & \texttt{ekta} & ektaa, akta, ekhta \\
now & \texttt{ekhon} & ekhone, akhon \\
is happening & \texttt{hocche} & hoche, hochhe, hoicche \\
is, exists & \texttt{ache} & achhe, ase \\
if & \texttt{jodi} & zodi \\
but & \texttt{kintu} & kintuh \\
then, so & \texttt{tahole} & tahle \\
good & \texttt{bhalo} & valo, bhaalo \\
plural marker & \texttt{gula} & gulo \\
\bottomrule
\end{tabular}
\end{table}

Technical terms stay in English and are spelled in English. Text is ASCII only, with no
smart quotes and no Bengali script.

\section*{A.5 The validator}

A script checks every transcript against the rules above and reports segments over 30
seconds, unreadable timestamps, gaps and overlaps, missing headers and spellings outside
the canonical list. It can repair malformed timestamps automatically, keeping a backup
copy of the original file. The checks that would corrupt training silently, such as a
transcript with no matching video or a lecturer label that disagrees with the voice
check, are a hard gate that stops a training run before it starts.
